\originalchapter{自由费米拓扑相与质量矩阵}
\label{ch:spt}
\source{8.513第16--19讲幻灯片；分类框架的数学定理按明确范围引用。}
\originalsection{有隙哈密顿量的稳定等价}
非相互作用、保持粒子数的多体哈密顿量$\hat H=\sum_{ab}M_{ab}c_a^\dagger c_b$由Hermitian单体矩阵决定。以费米能为零，有隙指$M$无零本征值。连续变形$M\to Q=\operatorname{sgn}M$保留本征矢与正负能级数，且$Q^2=1$。在零维，负能级数不能在不闭隙下改变；加上完全局域填满或空的平凡轨道后，分类使用稳定差值。

周期体系给每个$k$的占据投影$P(k)=(1-Q(k))/2$。若$P$有$p$维像与$q$维核，其值域为复Grassmann空间$U(p+q)/(U(p)\times U(q))$。高维拓扑分类不是只数某个$k$的负能级，而是研究整个Brillouin区到此空间的连续映射及其对称约束。加入平凡带把小矩阵的偶然障碍与稳定拓扑分开。
\originalsection{对称类与Clifford扩张}
时间反演$T$和粒子--空穴操作$C$为反幺正，分别满足$TH(k)T^{-1}=H(-k)$与$CH(k)C^{-1}=-H(-k)$；其平方可为$\pm1$。两者的乘积可给幺正手征对称$\Gamma H\Gamma^{-1}=-H$。由这些条件产生十种自由费米对称类。BdG中的粒子--空穴关系包含Nambu冗余，不应把它误当独立物理守恒电荷。

连续Dirac模型$H(k)=\sum_{i=1}^d k_i\gamma_i+mM$若要谱满足$E^2=k^2+m^2$，需$\{\gamma_i,\gamma_j\}=2\delta_{ij}$、$\{M,\gamma_i\}=0$、$M^2=1$。将对称操作在实Majorana基中转为与$M$对易或反对易的生成元，寻找允许质量相当于给Clifford代数增加一个生成元。允许质量空间的连通分支决定能否把两种符号质量连接而不闭隙。

Bott周期性给复分类空间两周期、实分类空间八周期。该同伦定理属于所调用的数学输入，本册不声称用有限几个Pauli矩阵例子证明其一般版本。应用时先明确自由、稳定、局域、有隙、保持指定内部对称的假设。常见强拓扑分类如下，晶格额外弱指标及晶体对称指标不含在表内。
\begin{center}\begin{tabular}{llll}\toprule 对称类&对称条件&空间维数&稳定强指标\\\midrule
A&仅粒子数守恒&2&$\mathbb Z$（Chern）\\
AIII&手征，无$T,C$&1&$\mathbb Z$（绕数）\\
D&$C^2=+1$&1&$\mathbb Z_2$\\
D&$C^2=+1$&2&$\mathbb Z$\\
AII&$T^2=-1$&2、3&$\mathbb Z_2$\\
BDI&$T^2=C^2=+1$&1&$\mathbb Z$\\\bottomrule\end{tabular}\end{center}
\originalsection{SSH链与Kitaev链}
SSH链的两个亚晶格跃迁$t_1,t_2$给$H(k)=(t_1+t_2\cos k)\sigma_x+t_2\sin k\sigma_y$。手征算符$\sigma_z$与$H$反对易。复函数$q(k)=t_1+t_2\e^{\ii k}$的闭曲线绕原点整数$\nu=(2\pi\ii)^{-1}\int q^{-1}\partial_kq\dd k$。$|t_2|>|t_1|$绕一圈，反之零；$|t_1|=|t_2|$闭隙。若加$\sigma_z$势差，手征破坏，可在不闭隙下绕开原点，这说明保护对称的作用。

Kitaev链$H=-\sum_j[tc_j^\dagger c_{j+1}+\Delta c_jc_{j+1}+\mathrm{h.c.}]-\mu\sum_j(n_j-1/2)$。取实$t,\Delta$，BdG谱$E_k^2=(-2t\cos k-\mu)^2+4\Delta^2\sin^2k$。在$\Delta\ne0$时闭隙只能于$k=0,\pi$且$\mu=\mp2t$。Majorana质量在这两点符号乘积为$\operatorname{sgn}[(\mu+2t)(\mu-2t)]$，负号区$|\mu|<2|t|$非平凡，开链有两端零模。允许保持$C$但破坏额外实时间反演的扰动时只剩$\mathbb Z_2$保护，不能用实参数特殊模型的整数分类替代一般D类。
\originalsection{体边对应与相互作用边界}
质量畴壁上的零模展示体边对应：允许质量由一相换到另一相，若无法全程保持有隙，界面必须出现低能态或其他相结构。边缘的具体速度、重构和附加平凡通道可变，体指标只固定净异常内容。

自由费米分类不是所有相互作用物态的分类。多体相可有内禀拓扑序，也可能使自由整数指标经相互作用约化；例如特定一维BDI体系在保持相应对称与局域作用下整数可约为八周期。这个例子只说明假设边界，本稿未推导所有相互作用分类。一般SPT是保持对称时不能连续化到平凡态的短程纠缠相，破坏保护对称后可平凡化；内禀拓扑序则不需这些保护对称，并有长程纠缠与任意子。
